\section{Fashioned for Presence}

The Ark begins with a command, not an inspiration in a craftsman's mind.
Israel is to make a sanctuary so that the Lord may dwell among his people, and
every part is to follow the pattern shown to Moses.  The first furnishing named
is a chest: two and a half cubits long, one and a half wide, one and a half
high; wood covered with pure gold within and without; four gold rings; poles
that remain in the rings; and above it a cover of pure gold beneath the wings
of two cherubim.  The testimony is placed within.  From above the cover,
between the cherubim, the Lord promises to meet Moses and speak his commands
for Israel (Exod 25:8--22).

Nothing in this description imprisons the God of heaven.  The same Scriptures
that locate the meeting also refuse containment.  The Ark is a real, appointed
place of covenant encounter precisely because God freely makes it so.  Its
holiness is received.  Its gold does not create glory; its dimensions do not
measure the divine nature.

\subsection{Wood, gold, and the Spirit-given craft}

The Hebrew names the timber \emph{shittim}, conventionally rendered acacia;
older English Catholic Bibles often retain ``setim'' wood.  The exact botanical
species is not certain enough to carry theology.  The ancient Greek Exodus
renders the phrase as ``incorruptible'' or non-decaying wood
(Exod 25:10 LXX).  That wording will later offer Christian contemplation a
fitting image of holy integrity.  It is not a claim that the material could
never decay, nor a direct prediction of a Marian dogma.  The literal object is
wood overlaid with gold, not a solid-gold chest (Exod 25:10--15; 37:1--5).

The work is at once human and graced.  Bezalel is called by name and filled
with the Spirit of God in wisdom, understanding, knowledge, and workmanship.
Oholiab and other skilled workers are given to the task.  Their art is neither
autonomous self-expression nor mechanical passivity: given intelligence
answers a revealed pattern.  Bezalel makes the Ark, its poles, the golden
cover, and the overshadowing cherubim; when the whole work is finished, Moses
sees that it has been done as the Lord commanded and blesses the workers
(Exod 31:1--11; 35:30--35; 37:1--9; 39:32, 42--43).

The union of hidden wood and visible gold invited later symbolic reading, but
the first lesson lies closer to the text.  God prepares a dwelling through the
free offerings of the people and the sanctified intelligence of artisans
(Exod 25:1--9).  Matter is not rejected.  It is ordered, offered, and made
servant to encounter.

\subsection{What stood in relation to the Ark}

The Bible's accounts of the sacred contents must be allowed to speak in their
own forms.  The tablets of the testimony are explicitly placed inside the Ark
(Exod 25:16, 21; 40:20; Deut 10:1--5).  The manna vessel is set before the
Lord or before the testimony (Exod 16:32--34).  Aaron's staff, after budding as
a sign against rebellion, is returned before the testimony (Num 17:1--11).
At the end of Deuteronomy the written book of the law is placed beside the Ark
as a witness, not inside it (Deut 31:24--26).

Hebrews later gathers the familiar threefold memory in its description of the
first covenant sanctuary: the golden manna vessel, Aaron's budding staff, and
the covenant tablets are associated with the Ark as things within it
(Heb 9:3--5).  At Solomon's dedication, however, Kings and Chronicles state
that only the two tablets are inside (1 Kgs 8:9; 2 Chr 5:10).  Scripture does
not narrate a timetable by which the manna or staff entered or left, and this
study will not invent one.  The differences themselves disclose a stable
center: the Ark gathers Israel's memory of God's word, provision, chosen
priesthood, and covenant testimony.

Deuteronomy's first-person recollection says that Moses made an Ark of wood and
placed the renewed tablets in it (Deut 10:1--5), while Exodus names Bezalel as
the craftsman of the Ark constructed according to the sanctuary pattern
(Exod 37:1).  The texts do not require an otherwise unattested second or
temporary Ark.  It is enough to preserve both modes of narration without
manufacturing the missing history.

\subsection{When glory filled the dwelling}

On the first day of the first month of the second year, Moses raises the
tabernacle.  He places the testimony in the Ark, sets the cover above it,
brings the Ark behind the veil, and finishes the work.  Then the cloud covers
the Tent of Meeting and the glory of the Lord fills the dwelling so completely
that Moses cannot enter (Exod 40:17, 20--21, 33--35).  What the artisans could
prepare, only God could fill.

The Greek text says that the cloud was ``overshadowing'' the Tent
(Exod 40:35 LXX).  The grammatical object is the Tent containing the installed
Ark, not the Ark by itself.  Yet the verb will sound again when Gabriel tells
Mary that the power of the Most High will overshadow her (Luke 1:35).  It is
the same verbal root, not the same inflected form and not a word unique to
these two passages.  Its force comes from the whole scene: a prepared
sanctuary, divine overshadowing, and the arrival of glory.

From then on, the cloud determines Israel's departures.  When it rises, the
people go forward; when it remains, they remain.  By day it is cloud, by night
fire, before the eyes of the whole house of Israel throughout their journeys
(Exod 40:36--38).  The newly fashioned Ark is already a traveling Ark.  The
glory that fills the sanctuary will now take the road.

\begin{studybox}{Preparation and indwelling}
The craftsmen supply no glory of their own.  They receive a pattern, offer
their skill, finish the commanded work, and wait upon God's coming.  This order
will govern the whole book: grace prepares a creature to receive a gift that
grace alone can give.  Christian typology may contemplate the prepared wood,
the gold, and the overshadowed sanctuary, but its center is always the Lord
who chooses to dwell.
\end{studybox}
